\section{A Deadline Model of Activist Stake Formation}
\label{sec:model}

The model provides the economic framework for the empirical evaluation. Its
main purpose is to derive the target-level exposure measure
$\mathrm{RTD}_i$ and to show when the statutory filing deadline limits
post-trigger accumulation. It also clarifies two downstream empirical
questions: whether earlier disclosure attracts or deters bidders, and whether
changes in measured takeover premia represent a reallocation of gains across
event time or a change in total target-shareholder value.

The distinctive margin is the time available to accumulate shares after
crossing the five-per-cent threshold. Existing disclosure models study a
stake-size threshold \parencite{ordonezcalafibernhardt2022}; they do not treat
the filing window as a deadline in an accumulation problem. Likewise,
\textcite{celentanolevine2025} provide the capitalisation interpretation of a
lower measured takeover premium, while \textcite{bett2014} provide the relevant
evidence on substitution between run-ups and markups. The model instead
connects the legal time limit to stake formation, bidder entry, and the
empirical decomposition of target gains.

\subsection{Setting and timing}
\label{sec:model-env}

A risk-neutral activist considers a campaign against target $i$, which has
$S_{\mathrm{out},i}$ shares outstanding and liquidity $L_i$. In the empirical
implementation, liquidity is measured by average daily volume over the six
months ending 60 trading days before the trigger, with Amihud illiquidity as an
alternative. The lag keeps the activist's untimed toehold purchases out of the
measurement window and aligns the liquidity measure with the unaffected-price
anchor. All agents discount at rate $\delta\in(0,1)$ per period.

At $t_0$, the activist privately observes an opportunity $\theta\geq0$, the
per-share value improvement available if the campaign succeeds. The activist
then builds a toehold and crosses the five-per-cent beneficial-ownership
threshold at $t_1$. This pre-trigger phase is not subject to the filing deadline
because the Section~13(d) clock begins only after the threshold is crossed
\parencite{sec2023modernization}.

The crossing date is the trigger date observed in the data. The unaffected
price $P^0$ is measured a fixed $k$ trading days before that date. This anchor
is necessary because accumulation itself can move the price. Slow toehold
building before trigger-minus-$k$ may still contaminate the measure, so the
empirical analysis varies $k$ across 30, 60, and 120 trading days.

Between the trigger and filing, the activist can add an undisclosed increment
$\Delta x$ over at most $d$ business days. The filing stake is
$x=0.05+\Delta x$. Before the reform, ten calendar days corresponded to
approximately seven business days; after the reform, $d=5$. The reform
therefore affects only the post-trigger increment.

Filing occurs at the earlier of the statutory deadline and voluntary
disclosure. Early filing remains feasible and may be optimal, so the deadline
binds only for an endogenous subset of campaigns. Let $P^F$ denote the
post-filing price and define $G_{\mathrm{fil}}\equiv\log(P^F/P^0)$. Potential
acquirers subsequently search and evaluate the target. Let $P^B$ denote the
bid-eve price and $G_{\mathrm{run}}\equiv\log(P^B/P^F)$. If bargaining produces
a transaction, the final offer is $B$ and
$G_{\mathrm{mkp}}\equiv\log(B/P^B)$; otherwise, the firm remains freestanding.

The activist's private information is summarised by the scalar $\theta$. The
model does not introduce a separate operational-versus-sale type. With a
two-dimensional type and a two-dimensional signal
$(x,\text{Item 4 purpose})$, the filing generically separates types and $P^F$
collapses to the full-information price, eliminating the signal noise needed
elsewhere in the model. Operational and sale intent therefore enter only as
empirical Item~4 heterogeneity in the entry elasticities below, not as a second
signalling dimension. Section~\ref{sec:model-omitted} discusses this restriction.

\subsection{The cost of accumulating before the deadline}
\label{sec:model-cost}

The activist trades against competitive dealers with linear price impact
$\lambda_L(L)$, where $\lambda_L$ falls with target liquidity. Empirically,
$\lambda_L$ is proportional to inverse average daily volume or to Amihud
illiquidity. Detection risk and dealer inventory capacity limit purchases to a
fraction $\phi$ of average daily volume per day.

At daily purchase rate $r$, excess execution cost is
$(\lambda_L/2)r^2$. Convexity makes a constant rate optimal over a fixed window.
The cost of acquiring $x$ over $d$ days is therefore
\begin{equation}
C(x;d,L)=\frac{\lambda_L(L)}{2}\frac{x^2}{d}.
\label{eq:cost}
\end{equation}
Its relevant derivatives are
\begin{equation}
C_x=\frac{\lambda_Lx}{d}>0,\qquad
C_{xx}=\frac{\lambda_L}{d}>0,\qquad
C_{xd}=-\frac{\lambda_Lx}{d^2}<0,
\label{eq:cost-derivs}
\end{equation}
and the horizon scaling is exact:
$C(x;\kappa d,L)=\kappa^{-1}C(x;d,L)$ for every $\kappa>0$.

A longer window lowers both total and marginal accumulation cost. The negative
cross-partial $C_{xd}$ is derived from the trading technology rather than
assumed. It is strongest for illiquid targets, for larger desired increments,
and close to the deadline. Moving from approximately seven to five business
days raises the quadratic cost of a fixed post-trigger increment by
$7/5=1.4$, or 40 per cent. The same price-impact mechanism produces a
pre-filing run-up proportional to $\lambda_Lx/d$, linking the split between
pre-filing leakage and the filing return to target liquidity.

Equation~\eqref{eq:cost} applies only to the post-trigger increment $\Delta x$.
The pre-trigger toehold is acquired over an unconstrained horizon chosen by the
activist, so its cost does not depend on $d$ and drops out of the comparative
statics.

This cost is a static reduced form for a dynamic trading problem. A deadline
version of the Kyle framework in \textcite{collindufresnefos2016} would make the
intensity path endogenous and allow the activist to trade off price impact
against information leakage. The strategic benchmark is \textcite{back2018}.
Those models add the shape of within-window trading, but not the sign of
$C_{xd}$ or the $1/d$ scaling whenever per-period impact is convex.
\textcite{collindufresnefos2015} establish that informed trading during the
Schedule~13D window is empirically important, which supports treating the
window as a trading horizon; \textcite{kyle1985} supplies the linear-impact
primitive.

\subsection{Exposure to the deadline and the stake at filing}
\label{sec:model-corner}

The daily participation limit implies that target $i$ can accommodate at most
\begin{equation}
\overline{\Delta x}_i(d)
=\frac{\phi\,\mathrm{ADV}_i\,d}{S_{\mathrm{out},i}}
\label{eq:cap}
\end{equation}
of undisclosed post-trigger accumulation. Define required trading days as
\begin{equation}
\mathrm{RTD}_i
=\frac{0.05\times S_{\mathrm{out},i}}
       {\phi\times\mathrm{ADV}_i},
\qquad\text{so that}\qquad
\overline{\Delta x}_i(d)=0.05\frac{d}{\mathrm{RTD}_i}.
\label{eq:rtd}
\end{equation}
This is the number of trading days required to add five percentage points of
ownership at the participation cap. Let
$x_i^\ast(d)=0.05+\Delta x_i^\ast(d)$ denote the stake at filing.

\begin{proposition}[Deadline exposure and stake formation]
\label{prop:R2}\label{prop:R3}
Let $\Delta x^u(d)$ solve the unconstrained post-trigger problem
\[
\max_{\Delta x\geq0}
\left\{b(0.05+\Delta x,d)-C(\Delta x;d,L)-F\right\},
\]
where $b$ is the gross benefit from the filing stake and $F$ is a fixed campaign
cost. The realised increment is
\[
\Delta x_i^\ast(d)
=\min\{\Delta x^u(d),\overline{\Delta x}_i(d)\}.
\]

First, the deadline does not constrain the activist's ability to reach the
disclosure threshold because the statutory clock starts at the crossing. It
constrains only ownership acquired beyond the threshold. The maximum filing
stake is therefore
\[
0.05\left(1+\frac{d}{\mathrm{RTD}_i}\right).
\]

Second, for a desired increment $\Delta x^u$, the constrained set is
\[
\left\{i:\mathrm{RTD}_i>\frac{0.05d}{\Delta x^u}\right\}.
\]
This set expands, in the set-inclusion sense, as $d$ falls. On the constrained
set, the shadow value
\[
b_x(0.05+\overline{\Delta x}_i(d),d)
-C_x(\overline{\Delta x}_i(d);d,L)\geq0
\]
increases with $\mathrm{RTD}_i$. Under the doubling calibration
$\Delta x^u=0.05$, a scale motivated by the pre-disclosure accumulation
documented by \textcite{collindufresnefos2015}, the constrained set is exactly
$\{i:\mathrm{RTD}_i>d\}$.

Third, suppose $b(\cdot,d)$ is $C^2$ and the objective is strictly concave in
$\Delta x$, so $b_{xx}-C_{xx}<0$, with derivatives of $b$ evaluated at
$x=0.05+\Delta x^\ast$. In the unconstrained region,
\begin{equation}
\frac{d\,\Delta x^\ast}{dd}
=-\frac{b_{xd}-C_{xd}}{b_{xx}-C_{xx}}>0
\qquad\Longleftrightarrow\qquad
b_{xd}>C_{xd}=-\frac{\lambda_L\Delta x^\ast}{d^2}.
\label{eq:sc}
\end{equation}
The filing stake moves one-for-one:
$dx^\ast/dd=d\,\Delta x^\ast/dd$. In the constrained region,
$\Delta x_i^\ast=\overline{\Delta x}_i(d)$ and
\[
\frac{d\,\Delta x_i^\ast}{dd}
=\frac{\phi\,\mathrm{ADV}_i}{S_{\mathrm{out},i}}>0
\]
without any additional restriction.

Consequently, the reform dose is a monotone function of
$\mathrm{RTD}_i/d$. The continuous dose is the primary model object. The
indicator $\mathbf{1}\{\mathrm{RTD}_i>d\}$ is the corresponding binary
classification under the doubling calibration, not a general treatment
definition.
\end{proposition}

Condition~\eqref{eq:sc} is the model's single-crossing restriction. A longer
window delays filing and hence certification, which may reduce the probability
that a bid arrives while the activist remains positioned. Thus $b_{xd}$ may be
negative. When it is, the condition requires
$|b_{xd}|<\lambda_L\Delta x^\ast/d^2$. The restriction is weakest where the
accumulation friction is largest: high $\lambda_L$, low $d$, and high
$\Delta x$. It is unnecessary at the constrained corner. The sign is therefore
conditional in the interior and unconditional only on the constrained set.

The proposition gives an economic interpretation to the accumulation ladder in
the adopting release \parencite{sec2023modernization}. Roughly 80 per cent of
filers complete their reported stake within five business days and belong to
the slack set for their realised desired increment. The 3 per cent of
non-corporate-action campaigns that remain below 90 per cent completion at the
amended deadline form the deeply constrained tail. Equations~\eqref{eq:cap} and
\eqref{eq:rtd} derive the exposure measure from lagged target-level, price-free
inputs. The measure is defined for one-time filers and permits
activist-by-quarter fixed effects.

\subsection{Bidder entry}
\label{sec:model-entry}

Potential acquirers are ex ante identical. Each pays search and evaluation cost
$s$ to learn its valuation, and entry continues until expected profit is zero.
Earlier disclosure has two opposing effects.

A Schedule~13D filing certifies that a large, sophisticated holder has
identified the target and is unlikely to defend incumbent management. The
filing therefore performs part of the bidder's screening and lowers $s$. This
is the solicitation and search mechanism of \textcite{corumlevit2019} and the
activist-intermediary mechanism of \textcite{burkartlee2022}, both taken as
given here.

Disclosure can also deter entry by increasing the expected acquisition price.
The disclosed stake and the associated engagement are capitalised into $P^F$,
and a larger, better-informed bloc may bargain more aggressively. This is the
pre-emption channel in \textcite{fishman1988} and the disclosure-reform evidence
in \textcite{bonettiduroormazabal2020}. It also contains the negative entry
effect in the companion model.

Let $h(s,B^e)$ be the free-entry bid hazard, decreasing in search cost and the
expected bid price. Let $e$ denote the earliness of disclosure, equal to minus
the filing lag, holding the campaign fixed. Define
\[
\varepsilon_{\mathrm{cert}}
=\left(-\frac{\partial\log h}{\partial\log s}\right)
 \left(-\frac{\partial\log s}{\partial e}\right),
\qquad
\varepsilon_{\mathrm{det}}
=\left(-\frac{\partial\log h}{\partial\log B^e}\right)
 \left(\frac{\partial\log B^e}{\partial e}\right).
\]
Both elasticities are non-negative, and
\begin{equation}
\frac{d\log h}{de}
=\varepsilon_{\mathrm{cert}}-\varepsilon_{\mathrm{det}}.
\label{eq:entry-elasticity}
\end{equation}
Earlier disclosure raises the bid hazard if and only if
$\varepsilon_{\mathrm{cert}}>\varepsilon_{\mathrm{det}}$. The model therefore
does not impose a sign on bidder entry.

Certification is more likely to dominate when screening is costly and the
counterparty set is unclear: opaque or hard-to-value assets, few natural
strategic buyers, cold industry M\&A markets, and Item~4 language indicating
sale or strategic-alternatives intent. Deterrence is more likely to dominate
when acquirers are already searching, as in hot merger markets or industries
with an obvious set of strategic buyers.

The reform does not change disclosure timing alone. By reducing $x^\ast$
through the preceding proposition, it can also lower $B^e$ and thereby increase
entry. This third effect partly offsets deterrence, so an empirical entry
response cannot be attributed entirely to certification.

\subsection{Bargaining and appropriability}
\label{sec:model-bargain}

Conditional on bidder entry, the final offer divides the surplus between the
acquirer's stand-alone valuation $V_a$ and the value $Y$ under a successful
campaign:
\begin{equation}
B=V_a+\beta(x)(Y-V_a),
\qquad
\beta(x)=\bar{\beta}\lambda_{\mathrm{app}},
\qquad
\lambda_{\mathrm{app}}\in[0,1].
\label{eq:bargain}
\end{equation}

The restriction on $\beta$ preserves the free-rider problem in
\textcite{grossmanhart1980freerider}. An unrestricted increasing $\beta(x)$
would assume that problem away. The analysis therefore imports only the
conclusion of the companion model's tender-game appendix:
\[
\lambda_{\mathrm{app}}=1-q(1-\gamma)\psi,
\]
where $q$ is the probability that a fringe raider appears at the disagreement
node, $\gamma\in[0,1]$ is the portability of the activist's improvement under a
different owner, and $\psi\in[0,1]$ measures pivotality. Appropriability falls
with fringe-raid risk and rises with portability.

Each primitive has an observable counterpart. Fringe M\&A intensity measured at
the industry-year level proxies for $q$. The mix of strategic and financial
acquirers, together with asset redeployability, proxies for $\gamma$. State and
charter antitakeover provisions index charter dilution, while the bloc's stake
relative to control thresholds indexes pivotality. Moreover,
$\lambda_{\mathrm{app}}$ rises discontinuously when the bloc crosses
the blocking threshold $1-\tau_c$. This produces a stake-size threshold
prediction linked directly to the deadline through $x$.

The tender game itself is not reproduced. Doing so would add an entry-cost
distribution, a synergy distribution, charter dilution, and a control threshold
as four additional primitives. To avoid a notation conflict, charter dilution,
denoted $\phi$ in the companion appendix, is denoted $\phi_c$ here; $\phi$
continues to denote the daily participation cap.

\subsection{Accounting for target gains}
\label{sec:model-measurement}

Total target gain between the unaffected price and the final offer decomposes
exactly as
\begin{equation}
G_{\mathrm{tot}}
\equiv\log\!\frac{B}{P^0}
=\underbrace{\log\!\frac{P^F}{P^0}}_{G_{\mathrm{fil}}}
+\underbrace{\log\!\frac{P^B}{P^F}}_{G_{\mathrm{run}}}
+\underbrace{\log\!\frac{B}{P^B}}_{G_{\mathrm{mkp}}}.
\label{eq:decomp}
\end{equation}
This is an accounting identity. Its empirical content comes from the
restrictions it places on changes in the three components.

Suppose the bid-eve price is the discounted expected resolution value,
\[
P^B=\delta\left[\varrho B+(1-\varrho)P^{\mathrm{fail}}\right],
\]
where $\varrho\in(0,1]$ is the completion probability and
$P^{\mathrm{fail}}<B$ is the freestanding value. Then
\[
G_{\mathrm{mkp}}
=\log(1/\delta)
-\log\!\left(\varrho+(1-\varrho)\frac{P^{\mathrm{fail}}}{B}\right)
\geq\log(1/\delta)>0.
\]
In simple returns, the measured markup is bounded below by
$(1-\delta)/\delta$ plus a strictly positive deal-risk term. A fully anticipated
deal therefore still has a positive measured markup; a declining markup alone
has no structural interpretation.

Holding the transaction fixed, so that $B$ and $P^0$ do not change,
equation~\eqref{eq:decomp} implies the testable restriction
\begin{equation}
dG_{\mathrm{mkp}}
=-\left(dG_{\mathrm{fil}}+dG_{\mathrm{run}}\right).
\label{eq:invariance}
\end{equation}
Under pure capitalisation, gains move across the three components while
$G_{\mathrm{tot}}$ remains invariant. If the reform destroys target value,
$G_{\mathrm{tot}}$ falls.

This restriction separates capitalisation from value destruction in the data.
\textcite{celentanolevine2025} give the capitalisation interpretation in their
structural counterfactual. \textcite{bett2014} document substitution between
value-adjusted run-ups and markups, whereas \textcite{schwert1996} find little
substitution without that adjustment. The empirical analysis therefore tests
the restriction on event-time-invariant components. A naive calendar-window
comparison would mechanically reassign days between the filing and run-up
components when the filing deadline changes, producing a truncation artefact
that runs opposite to the prediction being tested.

\subsection{Welfare and empirical sufficient statistics}
\label{sec:model-policy}

Let $W(d)$ denote expected target-shareholder gain. In the interior, the
envelope theorem removes the indirect effect of $d$ through $x^\ast$ from the
activist's own payoff. Any surviving $dx^\ast/dd$ term is therefore an
externality on target shareholders. At the constrained corner, the activist's
payoff changes at first order with the shadow value of the filing window, and
the stake response is
$dx^\ast/dd=\phi\,\mathrm{ADV}_i/S_{\mathrm{out}}$.

The welfare decomposition is
\begin{equation}
\begin{split}
\frac{dW}{dd}
&=\underbrace{\frac{\partial\Prob(\mathrm{bid})}{\partial d}
   \E\!\left[G_{\mathrm{tot}}\mid\text{marginal deal}\right]}
   _{\text{(i) entry margin}}
+\underbrace{\frac{\partial\Prob(\mathrm{camp})}{\partial d}
   \E\!\left[G_{\mathrm{tot}}\mid\text{marginal campaign}\right]}
   _{\text{(ii) campaign-supply margin}}\\[2pt]
&\qquad
+\underbrace{\Prob(\mathrm{bid})
   \frac{\partial\E[G_{\mathrm{mkp}}]}{\partial\beta}
   \beta'(x^\ast)\frac{dx^\ast}{dd}}
   _{\text{(iii) stake-bargaining margin}}.
\end{split}
\label{eq:welfare}
\end{equation}

The entry margin has the sign of equation~\eqref{eq:entry-elasticity} and must
be estimated. The campaign-supply margin is non-negative because a longer
window supports more campaigns. Since the reform has $\Delta d<0$, this term
enters the welfare effect negatively and captures chilling. The
stake-bargaining margin operates through the filing stake: a shorter window
lowers $x^\ast$, which lowers $\beta(x^\ast)$ and the conditional markup
independently of entry and campaign supply.

The conditional expectations in equation~\eqref{eq:welfare} concern marginal,
not average, units. Marginal deals are absent under the counterfactual regime,
and deterred campaigns are unobserved. Replacing these objects with the average
observed premium would substitute an ATE for a LATE. The intensity design
addresses this problem because campaigns with high $\mathrm{RTD}_i$ are the
units at the margin of being constrained or deterred. Its exposure-based
estimands therefore correspond to the weights required by the welfare
decomposition.

Finally, accumulation can move the target's price before filing. Part of
$G_{\mathrm{fil}}$ is consequently a transfer from selling shareholders to the
activist rather than a gain to a fixed shareholder cohort. Welfare is therefore
defined for shareholders present at trigger-minus-$k$, and the transfer
component is reported separately.

\subsection{Scope of the model}
\label{sec:model-omitted}

The filing price is a Bayesian update based on an exogenously garbled filing.
Nature adds noise to the disclosed $(x,\text{purpose})$ pair, and a false Item~4
statement carries an exogenous securities-liability cost. The activist's stake
is not solved as an endogenous signal. A continuous signal chosen by a privately
informed type can support pooling, partial-pooling, and separating equilibria.
Without an equilibrium refinement and a uniqueness result, $P^F$ would not be
a well-defined empirical object.

This restriction also avoids a separate closure problem. If the final offer
$B$ in equation~\eqref{eq:bargain} is independent of $P^F$, the capitalisation
result follows mechanically because the denominator moves while the numerator
does not. If $B$ instead depends on $P^F$ through a market-price floor, the
mapping $P^F\to B\to P^F$ becomes a fixed point and requires a contraction
condition. The model therefore holds $B$ independent of $P^F$ and presents the
invariance result as a measurement restriction rather than the
paper's main theoretical claim. Endogenising the filing price is left for
separate work.

Other deliberate omissions are the operational-versus-sale distinction as a
second private type, the tender game underlying $\lambda_{\mathrm{app}}$,
bidder competition and repeated fringe entry, and the within-window trading
path. Cash-settled derivatives are treated in the empirical Item~6 analysis as
an outside option but are not priced in the model. This omission is
consequential because derivative substitution is a primary empirical margin.
The model's narrower role is to identify the campaigns most likely to
substitute: under the doubling calibration, $\{\mathrm{RTD}_i>d\}$.

The resulting specification has one private type, quadratic accumulation cost,
an exogenously garbled filing, a free-entry bidder cutoff, and a bargaining
share disciplined by one appropriability parameter. Any additional mechanism
would require removing an existing one rather than expanding this core.
